% The editor's preface. DRAFT — Wilson's to edit; the voice should be his.
% Every claim here is grounded in the repo: CONVENTIONS.md for the rules,
% apparatus/THE-LATIN-QUESTION.md for printed 34/35 and 174/175,
% apparatus/CLASS-A-ledger.md for the dropped Hebrew, CLASS-B-unit-boundaries.md
% for the blank leaves after printed 184, EDITION-SHAPE.md for the layout.
% If you change a claim here, change it there too.
\begingroup
% The body of the book is unindented sense-lines, so the preface separates its
% paragraphs by space rather than by indent. Scoped: it must not leak into the text.
\setlength{\parskip}{0.55\baselineskip}
\setlength{\parindent}{0pt}
\addcontentsline{toc}{section}{Preface}
\markboth{Preface}{Preface}
{\centering\Large\scshape Preface\par}
\vspace{1.5\baselineskip}

Lancelot Andrewes kept a book of prayers in his own hand, in Greek, in Latin, and
in Hebrew, and did not publish it. What survives is a series of witnesses to it:
manuscripts that differ from one another, and printed editions that differ from
the manuscripts. This edition sets out one of those witnesses whole --- the text
printed at Oxford by J.\,H. Parker in 1853, itself a reprint of the Sheldonian
edition of 1675 --- and puts a new English translation beside it.

The English translations of Andrewes now in print all descend from one
nineteenth-century construal. No edition currently in print, and none in the
public domain, presents the Greek and the Latin at all. That is the gap this book
is for. A reader who wants to know what Andrewes actually wrote, rather than what
a translator made of it, has until now had to find a scan.

\vspace{0.8\baselineskip}
{\scshape The text.}\quad
What is printed here on the left-hand pages is a record of the 1853 plate, not an
improvement on it. Accents, breathings, punctuation, capitals and ligatures stand
as they were set. Where the 1853 gives a scripture reference that is wrong, the
wrong reference is printed and the note says so; where a sort is broken or a word
has been lost to a defect in the plate, the loss is marked rather than mended.
Nothing has been silently corrected. An edition that quietly repairs its source
is no longer a witness to anything, and this one is meant to be usable as
evidence.

The consequence is that a few things in this book are known to be wrong and are
left wrong on purpose. They are marked, and the reasons are given.

\vspace{0.8\baselineskip}
{\scshape Which language is Andrewes'.}\quad
In the first part the Greek is Andrewes' own text and the Latin faces it. The
translation here follows the Greek, and treats the Latin as the first witness to
how the Greek was construed by someone much nearer to it than we are. Where the
two genuinely diverge, the Greek is followed and the Latin's reading is recorded.
In the second and third parts, which are largely Latin, the Latin is the text and
is translated as such. Where Hebrew stands in the text it is kept, and translated;
the English layer never silently drops a script.

The 1853 editor was modest about his Latin. His apparatus says of the Greek-only
manuscript he collated that it lacks the Latin, \emph{quam Editor ipse confecisse
videtur} --- which the editor seems to have made himself. That would make the
Latin of the first half of Part I an editor's Latin, and worth less as a witness
than it looks.

It is worth more than he claimed. At printed 34--35 the manuscript reads
πλάσμα where the printed Greek reads τὰ ἔργα, and the facing Latin reads
\emph{Figmenta} --- which departs from the printed Greek, and departs also from
the Vulgate of the psalm it is quoting, which reads \emph{opera manuum tuarum}.
There is no route to \emph{Figmenta} except from a Greek text reading πλάσμα. At
printed 174--175 the Latin \emph{Ablue} renders ἀπόλουσαι, wash thou away, where
the printed Greek has ἀπόλυσαι, release thou --- one letter apart, and the Latin
saw the other one.
The Latin column is therefore not simply downstream of the Greek beside it: in
places it sees a manuscript the printed Greek does not. That is a small finding
with a large consequence for how this book should be read, and it is why the
Latin is set as a column and not as a footnote.

\vspace{0.8\baselineskip}
{\scshape The English.}\quad
The translation is new, and it is line-keyed: one English line to one line of the
original, at the same indentation, with the same brace-groups and columns. A
reader should be able to run a finger down the page and stay in step. That
constraint costs something in fluency and is worth its cost, because the shape of
these prayers on the page --- the lists, the widening and narrowing, the single
words given a line to themselves --- is not decoration. It is how Andrewes prayed.

The register is the English of the Authorised Version and the Prayer Book:
\emph{thou} and \emph{thee} to God, and the cadences an English reader already
carries. The syntax, though, stays plain. Archaism is confined to the pronouns and
verb-endings of address; there are no Jacobean inversions the Greek does not
force. For the Psalms the crib is Coverdale's Psalter, the Prayer Book's, ahead of
the Authorised Version --- not from preference but because Coverdale worked
through the Latin and the Greek, as Andrewes did, and so agrees with him in
exactly the places where the Authorised Version, made from the Hebrew, does not.
Andrewes never quoted an English Bible. He was quoting the Septuagint and the
Vulgate, and Coverdale stands nearer to that than any later English does.

\vspace{0.8\baselineskip}
{\scshape The page.}\quad
The originals are on the left-hand page and the English on the right. In Part I,
where the 1853 sets Greek and Latin on facing pages, both share the left-hand page
in two columns, so that the whole of an opening of the 1853 stands under the eye
at once with its translation beside it.

An earlier plan reproduced the 1675 opening exactly --- Greek verso, Latin recto,
page for page, with the English as a register along the foot. It was the more
beautiful arrangement and it was given up, for a reason worth recording: it made
every page of the book depend on every page fitting. It also had no room for the
English in those parts of the book that are in one language only, where there is
no facing page to spend. The present arrangement never depends on a page fitting,
and gives the English a page of its own throughout.

The page numbers running along the outer edge are this edition's. The numbers in
brackets on the inner edge are the 1853's, and every page of it is here and
findable by them.

\vspace{0.8\baselineskip}
{\scshape What is not here.}\quad
Two things, and it is better to say so.

The 1853 does not print all of Andrewes' Hebrew. Its own apparatus records Hebrew
standing in the manuscript at more than a dozen places where the printed page has
none --- among them a phrase of rabbinic anthropology, the \emph{tohu wa-bohu} of
the second verse of Genesis, and a short Greek-and-Hebrew lexicon of the words for
sin. This edition prints the 1853, so that Hebrew is recorded in the apparatus
and is not restored to the text. Restoring it would be a different book: a
critical text of Andrewes rather than a faithful setting of a printed edition, and
it would need the manuscripts themselves.

And the manuscript has silences the print does not show. Three blank leaves stand
in it after what is here printed page 184, in the middle of the seventh day. What
that silence means --- whether a division has been lost, or something was meant to
be added --- cannot be settled from a printed edition. It is recorded and left
open.

\vspace{0.8\baselineskip}
{\scshape Acknowledgement.}\quad
F.\,E. Brightman's translation of 1903 remains the best English understanding of
this book, and has been consulted throughout --- after drafting, never before, and
never copied. His arrangement is his own: he judged the second part of the 1675
unusable as printed and rearranged the whole for practical use. That judgement is
defensible, and its consequence is that the \emph{Preces Privatae} as a book, in
its own order, has been out of reach of English readers for a long time. Restoring
that order is much of what this edition is for. Where his notes have corrected a
reference or settled an attribution here, the note says so.

\vspace{1.2\baselineskip}
{\raggedleft\itshape Wroot Press\par}
\endgroup
\clearpage
